\documentclass[12pt,a4paper]{article}
\usepackage[utf8]{inputenc}
\usepackage[T1]{fontenc}
\usepackage[indonesian]{babel}
\usepackage{geometry}
\usepackage{longtable}
\usepackage{array}
\usepackage{booktabs}
\usepackage{enumitem}
\usepackage[hidelinks,breaklinks]{hyperref}
\geometry{margin=2.5cm}

\begin{document}

\begin{center}
\textbf{\Large RENCANA PEMBELAJARAN SEMESTER (RPS)}\\
\textbf{Berbasis Outcome-Based Education (OBE)}\\[0.5cm]
\textbf{PROGRAM STUDI TEKNIK INFORMATIKA}\\
\textbf{FAKULTAS TEKNOLOGI INFORMASI}\\
\textbf{UNIVERSITAS BALE BANDUNG}\\
\end{center}

\vspace{1cm}

% ============================================================
% 1. IDENTITAS MATA KULIAH
% ============================================================
\section*{1. Identitas Mata Kuliah}
\begin{tabular}{ll}
Nama Program Studi & : S1 Teknik Informatika \\
Nama Mata Kuliah & : Logika Matematika \\
Kode Mata Kuliah & : IT-231 \\ % Assuming a plausible code
Semester & : 2 (Dua) \\ % Assuming a plausible semester
SKS / Bobot Kredit & : 2 SKS \\
Dosen Pengampu & : (Nama Dosen Pengampu) \\
Tanggal Penyusunan & : \today \\
\end{tabular}

\vspace{0.5cm}

% ============================================================
% 2. CAPAIAN PEMBELAJARAN LULUSAN (CPL)
% ============================================================
\section*{2. Capaian Pembelajaran Lulusan (CPL)}

CPL yang dibebankan pada mata kuliah ini mencakup kompetensi lulusan dalam aspek pengetahuan, keterampilan, dan sikap:

\begin{itemize}[leftmargin=*]
  \item \textbf{CPL01:} Kemampuan menganalisis persoalan computing yang kompleks serta menerapkan prinsip-prinsip computing dan disiplin ilmu relevan lainnya untuk mengidentifikasi solusi, dengan mempertimbangkan wawasan perkembangan ilmu transdisiplin.
  \item \textbf{CPL02:} Memiliki pengetahuan yang memadai terkait dengan cara kerja sistem komputer dan mampu merancang dan mengembangkan berbagai algoritma metode untuk memecahkan masalah.
\end{itemize}

\vspace{0.5cm}

% ============================================================
% 3. CAPAIAN PEMBELAJARAN MATA KULIAH (CPMK)
% ============================================================
\section*{3. Capaian Pembelajaran Mata Kuliah (CPMK)}

Kemampuan atau kompetensi spesifik yang diharapkan mahasiswa kuasai setelah menyelesaikan mata kuliah:

\begin{itemize}[leftmargin=*]
  \item \textbf{CPMK1 (Fokus Logika Proposisional \& Analisis Argumen - Terkait CPL01):} Mahasiswa mampu menganalisis validitas argumen dan persoalan komputasi dasar dengan menerapkan prinsip-prinsip logika proposisional, tabel kebenaran, dan penarikan kesimpulan (silogisme, modus ponens, modus tollens) untuk mengidentifikasi solusi yang logis.
  \item \textbf{CPMK2 (Fokus Logika Predikat \& Pemodelan - Terkait CPL01 \& CPL02):} Mahasiswa mampu memodelkan persoalan yang lebih kompleks ke dalam bentuk ekspresi logika predikat (menggunakan kuantor universal dan eksistensial) sebagai langkah awal dalam merancang metode penyelesaian masalah yang terstruktur.
  \item \textbf{CPMK3 (Fokus Aljabar Boolean \& Sistem Komputer - Terkait CPL02):} Mahasiswa mampu menerapkan hukum-hukum aljabar Boolean untuk menyederhanakan fungsi logika dan merepresentasikannya ke dalam bentuk gerbang logika dasar, sebagai bentuk pemahaman terhadap cara kerja perangkat keras sistem komputer.
  \item \textbf{CPMK4 (Fokus Teknik Pembuktian \& Validasi Algoritma - Terkait CPL01 \& CPL02):} Mahasiswa mampu merancang argumen matematis dan membuktikan kebenaran suatu metode atau algoritma sederhana menggunakan teknik pembuktian formal (seperti deduksi, kontradiksi, atau induksi matematika) untuk memastikan solusi komputasi berjalan sesuai spesifikasi.
\end{itemize}

\vspace{0.5cm}

% ============================================================
% 4. SUB-CPMK / INDIKATOR PENCAPAIAN
% ============================================================
\section*{4. Sub-CPMK / Indikator Pencapaian}

Penjabaran CPMK menjadi indikator yang lebih terukur dan dapat diuji:

\begin{itemize}[leftmargin=*]
  \item \textbf{Sub-CPMK 1.1:} Mahasiswa mampu menjelaskan konsep dasar proposisi, nilai kebenaran, dan kata hubung logika (konjungsi, disjungsi, negasi, implikasi, biimplikasi).
  \item \textbf{Sub-CPMK 1.2:} Mahasiswa mampu menggunakan tabel kebenaran untuk mengevaluasi pernyataan majemuk serta mengidentifikasi tautologi, kontradiksi, dan kontingensi.
  \item \textbf{Sub-CPMK 1.3:} Mahasiswa mampu menganalisis validitas argumen menggunakan aturan inferensi (modus ponens, modus tollens, silogisme) dalam memecahkan persoalan komputasi dasar.
  \item \textbf{Sub-CPMK 2.1:} Mahasiswa mampu menjelaskan konsep dasar logika predikat, fungsi proposisional, dan ruang lingkup pembicaraan (\textit{universe of discourse}).
  \item \textbf{Sub-CPMK 2.2:} Mahasiswa mampu menerjemahkan pernyataan bahasa alami ke dalam ekspresi logika predikat menggunakan kuantor universal dan eksistensial.
  \item \textbf{Sub-CPMK 2.3:} Mahasiswa mampu memodelkan masalah komputasi yang terstruktur dengan merepresentasikan logika predikat sebagai basis perancangan algoritma.
  \item \textbf{Sub-CPMK 3.1:} Mahasiswa mampu mendefinisikan prinsip-prinsip dan hukum-hukum dasar Aljabar Boolean.
  \item \textbf{Sub-CPMK 3.2:} Mahasiswa mampu menyederhanakan fungsi logika menggunakan hukum Aljabar Boolean atau Peta Karnaugh (K-Map).
  \item \textbf{Sub-CPMK 3.3:} Mahasiswa mampu mengonversi dan merepresentasikan fungsi logika yang telah disederhanakan ke dalam bentuk rangkaian gerbang logika dasar pembentuk sistem komputer.
  \item \textbf{Sub-CPMK 4.1:} Mahasiswa mampu menerapkan berbagai metode pembuktian matematis dasar (pembuktian langsung, pembuktian tidak langsung/kontradiksi, dan kontraposisi).
  \item \textbf{Sub-CPMK 4.2:} Mahasiswa mampu mengaplikasikan prinsip Induksi Matematika untuk membuktikan kebenaran sebuah deret atau properti dari algoritma rekursif sederhana.
  \item \textbf{Sub-CPMK 4.3:} Mahasiswa mampu menyusun argumen matematis formal untuk memvalidasi langkah-langkah dalam penyelesaian masalah komputasi.
\end{itemize}

\vspace{0.5cm}

% ============================================================
% 5. MATERI PEMBELAJARAN (BAHAN KAJIAN)
% ============================================================
\section*{5. Materi Pembelajaran (Bahan Kajian) - Pemetaan 16 Pertemuan}

Daftar topik materi yang relevan dengan Sub-CPMK dan CPMK yang didistribusikan dalam 16 pertemuan (14 tatap muka, 1 UTS, 1 UAS):

\begin{enumerate}[leftmargin=*]
  \item \textbf{Pertemuan 1:} Pengantar Logika Matematika, Konsep Dasar Proposisi, dan Nilai Kebenaran (Sub-CPMK 1.1)
  \item \textbf{Pertemuan 2:} Kata Hubung Logika: Konjungsi, Disjungsi, Negasi, Implikasi, Biimplikasi (Sub-CPMK 1.1)
  \item \textbf{Pertemuan 3:} Evaluasi Pernyataan Majemuk dengan Tabel Kebenaran (Sub-CPMK 1.2)
  \item \textbf{Pertemuan 4:} Tautologi, Kontradiksi, Ekivalensi Logis, dan Kontingensi (Sub-CPMK 1.2)
  \item \textbf{Pertemuan 5:} Aturan Inferensi dan Penarikan Kesimpulan: Modus Ponens, Modus Tollens, Silogisme (Sub-CPMK 1.3)
  \item \textbf{Pertemuan 6:} Analisis Validitas Argumen dalam Konteks Komputasi Dasar (Sub-CPMK 1.3)
  \item \textbf{Pertemuan 7:} Pengantar Logika Predikat, Fungsi Proposisional, dan \textit{Universe of Discourse} (Sub-CPMK 2.1)
  \item \textbf{Pertemuan 8:} Ujian Tengah Semester (UTS)
  \item \textbf{Pertemuan 9:} Kuantor Universal dan Kuantor Eksistensial, Terjemahan Bahasa Alami ke Logika Predikat (Sub-CPMK 2.2)
  \item \textbf{Pertemuan 10:} Pemodelan Masalah Komputasi menggunakan Logika Predikat (Sub-CPMK 2.3)
  \item \textbf{Pertemuan 11:} Prinsip-prinsip dan Hukum-hukum Dasar Aljabar Boolean (Sub-CPMK 3.1)
  \item \textbf{Pertemuan 12:} Penyederhanaan Fungsi Logika menggunakan Hukum Aljabar Boolean dan Peta Karnaugh (K-Map) (Sub-CPMK 3.2)
  \item \textbf{Pertemuan 13:} Representasi Fungsi Logika ke dalam Rangkaian Gerbang Logika Dasar (Sub-CPMK 3.3)
  \item \textbf{Pertemuan 14:} Metode Pembuktian Matematis: Pembuktian Langsung, Kontradiksi, Kontraposisi (Sub-CPMK 4.1)
  \item \textbf{Pertemuan 15:} Induksi Matematika dan Penerapannya pada Algoritma Rekursif (Sub-CPMK 4.2, 4.3)
  \item \textbf{Pertemuan 16:} Ujian Akhir Semester (UAS)
\end{enumerate}

\vspace{0.5cm}

% ============================================================
% 6. METODE PEMBELAJARAN
% ============================================================
\section*{6. Metode Pembelajaran}

Strategi atau pendekatan pembelajaran yang dipilih sesuai OBE yang menekankan aktivitas mahasiswa:

\begin{itemize}[leftmargin=*]
  \item \textbf{Ceramah Interaktif:} Penjelasan konsep-konsep dasar logika dengan sesi tanya jawab aktif.
  \item \textbf{Problem-Based Learning (PBL):} Pemberian studi kasus terkait komputasi dasar untuk dianalisis dan diselesaikan menggunakan prinsip logika dan sistem inferensi.
  \item \textbf{Diskusi Kelompok:} Mahasiswa berdiskusi untuk menerjemahkan spesifikasi sistem bahasa alami menjadi model logika predikat yang ketat.
  \item \textbf{Praktikum/Latihan Mandiri:} Latihan terstruktur menyederhanakan rangkaian logika menggunakan Aljabar Boolean dan K-Map, serta merancang gerbang logika dasarnya.
  \item \textbf{Review Teman Sejawat (Peer Review):} Saling memeriksa argumen pembuktian matematis antara mahasiswa untuk mengasah pemikiran kritis terhadap kelemahan argumen.
\end{itemize}

\vspace{0.5cm}

% ============================================================
% 7. PENGALAMAN BELAJAR MAHASISWA
% ============================================================
\section*{7. Pengalaman Belajar Mahasiswa}

Deskripsi tugas, aktivitas, atau pengalaman belajar yang mendukung pencapaian Sub-CPMK:

\begin{itemize}[leftmargin=*]
  \item Menyusun dan mengevaluasi status kebenaran dari pernyataan logika majemuk yang sering ditemukan dalam percabangan algoritma (If-Else statement).
  \item Melakukan analisis validitas terhadap pernyataan (argumen) yang mendasari proses pengambilan keputusan (Decision Making) dalam suatu sistem komputasi kognitif sederhana.
  \item Merancang spesifikasi model sistem perangkat lunak skala kecil yang direpresentasikan oleh formula logika predikat menggunakan berbagai jenis kuantor.
  \item Terlibat aktif dalam mereduksi atau mengoptimalkan ekspresi boolean yang rumit sehingga mencapai jumlah gerbang logika aritmatika minimum yang paling efisien, sebelum mengimplementasikannya secara fisis/simulasi.
  \item Menulis dan mempresentasikan analisis pembuktian algoritma sederhana (contoh perhitungan iteratif atau rekursif faktorial) yang telah mereka buat secara matematis (dengan teknik induksi).
\end{itemize}

\vspace{0.5cm}

% ============================================================
% 8. KRITERIA, INDIKATOR, DAN BOBOT PENILAIAN
% ============================================================
\section*{8. Kriteria, Indikator, dan Bobot Penilaian}

Teknik/alat asesmen dipetakan ke Sub-CPMK/CPMK dengan bobot yang jelas:

\begin{longtable}{|>{\raggedright\arraybackslash}p{2.5cm}|>{\raggedright\arraybackslash}p{4cm}|>{\raggedright\arraybackslash}p{5.5cm}|c|}
\hline
\textbf{Komponen} & \textbf{Teknik Asesmen} & \textbf{Indikator/CPMK} & \textbf{Bobot (\%)} \\
\hline
\endfirsthead
\hline
\textbf{Komponen} & \textbf{Teknik Asesmen} & \textbf{Indikator/CPMK} & \textbf{Bobot (\%)} \\
\hline
\endhead
\hline
\endfoot

Tugas Individu & Latihan Soal Tertulis \& Pemodelan Khusus & Sub-CPMK 1.1, 1.2, 2.1, 2.2, 3.1 & 20 \\
\hline
Kuis & Pilihan Ganda \& Esai Pendek & Sub-CPMK 1.3, 3.2 & 10 \\
\hline
Tugas Kelompok & Proyek Mini (Pembuktian Algoritma \& Desain Fungsi) & Sub-CPMK 2.3, 3.3, 4.3 & 20 \\
\hline
UTS & Ujian Tulis (Kasus Komprehensif) & CPMK-1, Sebagian CPMK-2 & 25 \\
\hline
UAS & Ujian Tulis (Kasus Komprehensif) & CPMK-2, CPMK-3, CPMK-4 & 25 \\
\hline
\textbf{Total} & & & \textbf{100} \\
\hline
\end{longtable}

\textbf{Kriteria Penilaian:}
\begin{itemize}[leftmargin=*]
  \item A (85-100): Menguasai semua CPMK dengan sangat baik, argumen logis sempurna.
  \item B (70-84): Menguasai sebagian besar CPMK dengan baik.
  \item C (60-69): Menguasai CPMK dasar dengan cukup.
  \item D (50-59): Menguasai sebagian kecil CPMK.
  \item E ($<$50): Belum menguasai CPMK yang ditetapkan.
\end{itemize}

\vspace{0.5cm}

% ============================================================
% 9. EVALUASI DAN REFLEKSI PEMBELAJARAN (OPSIONAL)
% ============================================================
\section*{9. Evaluasi dan Refleksi Pembelajaran}

Penilaian sumatif/formatif untuk memantau ketercapaian outcome secara menyeluruh:

\begin{itemize}[leftmargin=*]
  \item \textbf{Evaluasi Formatif:} Analisis jawaban kuis mingguan dan respons dari peer review untuk memberikan intervensi belajar secara cepat jika mayoritas mahasiswa mengalami kesulitan di logika tertentu.
  \item \textbf{Evaluasi Sumatif:} Hasil UTS dan UAS akan dipetakan langsung ke pencapaian tiap CPL untuk menjustifikasi tercapai atau belumnya profil lulusan bagian dasar komputasi logis.
  \item \textbf{Continuous Improvement:} Kesalahan logika (logical fallacy) atau argumen yang paling tidak valid saat evaluasi tahun ini, akan menjadi titik fokus latihan studi kasus terbanyak untuk RPS di tahun ajaran depan.
\end{itemize}

\vspace{0.5cm}

% ============================================================
% 10. DAFTAR REFERENSI
% ============================================================
\section*{10. Daftar Referensi}

Sumber belajar utama yang digunakan dalam penyusunan materi dan asesmen:

\begin{enumerate}[leftmargin=*]
  \item The Open Logic Project. \textit{The Open Logic Text}. [OER - Buku teks pengenalan formal logika, Set Theory, and Proofs].
  \item Gallier, Jean H. \textit{Logic for Computer Science: Foundations of Automatic Theorem Proving}. [Catatan terbuka, University of Pennsylvania].
  \item Hammack, Richard. \textit{Book of Proof}. [OER - Buku yang fokus dalam mengenalkan dan menajamkan teknik bukti matematis].
  \item Modul Logika Informatika / Bahan Ajar Logika Matematika dari berbagai repositori institusional di Indonesia (UI, ITB, Unesa, dll).
  \item Seri Kursus Terbuka (OCW): MIT OCW 6.042J \textit{Mathematics for Computer Science}.
\end{enumerate}

\vspace{1cm}

\begin{flushright}
\begin{tabular}{c}
Bandung, \today\\[2cm]
\textbf{Nama Dosen Pengampu}\\
NIDN. XXXXXXXXXX
\end{tabular}
\end{flushright}

\end{document}
